\subsection{The Discovery Cycle: {{PERCENT_TRAITS_SATURATED}}\% of Traits Show GWAS Saturation}

\begin{figure}[t!]
    \centering
    {{fig:discovery_cycle}}
    \caption{The genetic discovery cycle across the phenome. (a) Schematic illustrating the expected discovery metric and saturation detection approach. (b-e) Exemplar traits representing different discovery stages: (b) Parkinson's disease (early, power-limited---all curves linear), (c) Asthma (mid-stage, annotation-limited), (d) Type 2 diabetes (saturated, pathway-concentrated---indirect diverges from direct), (e) Height (saturated, distributed---all curves saturate together). (f) Distribution of discovery status across {{NUM_TRAITS_SATURATION_ANALYSIS}} common traits.}
    \label{fig:discovery-cycle}
\end{figure}

To understand where research resources should be allocated across the phenome, we developed a framework for diagnosing each trait's position in the genetic discovery cycle. For any given trait, GWAS discovery follows a characteristic trajectory: early stages are limited by statistical power, but eventually the rate of new gene discoveries saturates even as sample sizes grow. Distinguishing power-limited from saturated traits is critical for research prioritization---continued GWAS investment yields diminishing returns for saturated traits, whereas power-limited traits benefit from larger cohorts.

We measure gene discovery using an \textit{expected discovery} metric that quantifies how much evidence has accumulated for each gene above a baseline prior. For any evidence type (direct, indirect, or combined), the expected discovery for gene $i$ is the posterior probability of disease relevance minus the prior probability: $D_i = P_i^{\text{post}} - P_i^{\text{prior}}$, where $P_i^{\text{post}} = \pi \cdot \text{ABF}_i / (1 + \pi \cdot \text{ABF}_i)$ and $\pi = 0.05$ is the background prior (\ref{sec:methods}). The total expected discovery $\tilde{N} = \sum_i D_i$ represents the expected number of genes identified above baseline---a continuous metric that avoids arbitrary significance thresholds while remaining interpretable as a gene count. Results were robust to the choice of background prior (Supplementary Figure~\ref{fig:prior-sensitivity}).

To trace how gene discovery accumulates with statistical power, we applied a downsampling procedure to {{NUM_TRAITS_SATURATION_ANALYSIS}} common traits with meta-analyzed GWAS summary statistics. For each trait, we simulated reduced sample sizes by inflating the standard error of each SNP's effect estimate according to the downsampled fraction $f \in \{0.1, 0.2, \ldots, 1.0\}$ (\ref{sec:methods}). At each fraction, we computed expected discovery separately for direct support ($\tilde{N}_D$), indirect support ($\tilde{N}_I$), and combined support ($\tilde{N}_C$), yielding discovery trajectories as a function of effective sample size.

We classified each trait's discovery status by comparing the fit of two models to its trajectory. A linear model implies that gene discovery continues to increase proportionally with sample size---the trait remains power-limited. An exponential saturation model implies that discovery is approaching an asymptote---the trait has reached or is entering saturation. We compared models using the Bayesian Information Criterion and further quantified the degree of saturation as the proximity to the fitted asymptote (\ref{sec:methods}).

Across {{NUM_TRAITS_SATURATION_ANALYSIS}} common traits, {{PERCENT_TRAITS_SATURATED}}\% showed evidence of GWAS saturation for direct genetic support (Figure~\ref{fig:discovery-cycle}f). The remaining traits split between power-limited ({{NUM_TRAITS_POWER_LIMITED}} traits) and annotation-limited ({{NUM_TRAITS_ANNOTATION_LIMITED}} traits). Among saturated traits, we observed two distinct architectural patterns based on the relationship between direct and indirect discovery trajectories.

Parkinson's disease exemplifies a trait in the early discovery stage (Figure~\ref{fig:discovery-cycle}b). All discovery curves---direct, indirect, combined, and the number of independent loci---increase linearly with sample size, with no evidence of saturation. The parallel trajectories indicate that both disease-associated and disease-relevant genes are being discovered at comparable rates, and continued GWAS investment will yield new loci that in turn enable identification of additional relevant genes through indirect support.

Asthma represents a mid-stage trait where direct discovery is progressing but indirect support lags (Figure~\ref{fig:discovery-cycle}c). Direct support shows early signs of curvature while the number of independent loci continues to grow linearly. Critically, indirect support accumulates more slowly than direct support, suggesting that current functional annotations incompletely capture asthma-relevant biology. For such traits, investment in disease-relevant annotations---airway epithelial expression programs or immune pathway curation---may yield greater returns than additional GWAS sample size alone.

Type 2 diabetes exemplifies a saturated trait with pathway-concentrated genetic architecture (Figure~\ref{fig:discovery-cycle}d). Both direct and indirect support show clear exponential saturation, reaching their asymptotes by approximately 40\% of the full sample size. However, combined support diverges upward from direct support, and the number of independent loci continues increasing linearly even as per-gene discovery saturates. This pattern indicates that T2D biology is organized around discrete pathways: once anchored by associated genes, these pathways propagate indirect support to functionally related genes that lack direct associations. For T2D, the discovery frontier lies in understanding pathway relationships rather than finding additional GWAS loci.

Height presents a contrasting saturation pattern (Figure~\ref{fig:discovery-cycle}e). Like T2D, both direct and indirect support saturate early---by approximately 30\% of the full sample size. Unlike T2D, all curves saturate together without divergence: direct, indirect, and combined support reach similar asymptotes. This pattern reflects height's highly polygenic, distributed genetic architecture \cite{yengo_saturated_2022}: thousands of genes contribute small effects through diverse biological processes, and shared functional annotations do not strongly connect associated genes to additional relevant genes beyond those already detected.

These results provide a principled basis for allocating resources across the phenome. For the {{NUM_TRAITS_POWER_LIMITED}} power-limited traits, continued investment in larger GWAS cohorts is justified. For the {{NUM_TRAITS_ANNOTATION_LIMITED}} annotation-limited traits, resources should be directed toward generating disease-relevant functional annotations. For the {{PERCENT_TRAITS_SATURATED}}\% of saturated traits, the discovery frontier has shifted from finding new associations to mechanistic understanding, and for pathway-concentrated traits specifically, indirect support provides a route to disease-relevant genes beyond the reach of GWAS alone.

% ANALYSIS IDEA: Temporal validation - traits classified as "power-limited" in older GWAS
% should show more new discoveries in newer GWAS than traits classified as "saturated."

% FIGURE SUGGESTION: Supplementary figure showing all ~1000 traits in 2D space
% (saturation degree vs. direct/indirect divergence), colored by disease category.
